\documentclass[notitlepage,onecolumn,showpacs,aps,prb]{revtex4-2}
\usepackage{geometry}
\usepackage[bottom]{footmisc}
\usepackage{color}
%\usepackage{doi}
\usepackage{hyperref}
%\hypersetup{hidelinks=true,}
%\usepackage[T1]{fontenc}
%\usepackage[utf8]{inputenc}
\usepackage{amssymb}
%\usepackage[nointegrals]{wasysym}
\usepackage{graphicx}
\usepackage{amsmath}
%\usepackage[breakable]{tcolorbox}
%\usepackage{esint}
\usepackage{epstopdf}
%\usepackage{color}
\usepackage{caption}
\usepackage{subcaption}
\usepackage{tensor}
\usepackage{slashed}
\usepackage{cancel}
%\usepackage{sidecap}
%\sidecaptionvpos{figure}{t}
\allowdisplaybreaks


\usepackage{feynmp-auto}
\setlength{\unitlength}{1mm}
\DeclareGraphicsRule{*}{mps}{*}{}

\newcommand{\bsym}[1]{\boldsymbol{#1}}
\newcommand{\mbf}[1]{\mathbf{#1}}
\newcommand{\op}[1]{\hat{#1}}
\newcommand{\grad}{\boldsymbol{\nabla}}
\newcommand{\Tr}{\text{Tr}}
\newcommand{\tr}{\text{tr}}
\newcommand{\N}{\mathcal{N}}
\renewcommand{\k}{\mathbf{k}}
\newcommand{\q}{\mathbf{q}}
\renewcommand{\r}{\mathbf{r}}
\renewcommand{\d}{\mathfrak{D}}
\newcommand{\F}{\mathcal{F}}
\renewcommand{\L}{\mathcal{L}}
%\newcommand{\beq}{\begin{equation}\begin{aligned}}
%\newcommand{\eeq}{\end{aligned}\end{equation}}

%\newcommand{\beq}{\begin{align}}
%\newcommand{\eeq}{\end{align}}

\def\beq#1\eeq{\begin{align}#1\end{align}}

\newcommand{\todo}[1]{{\color{red} {\bf [ToDo: #1]}}}
\newcommand{\mv}[1]{{\color{green}{#1}}}

\DeclareMathOperator*{\sumint}{%
\mathchoice%
  {\ooalign{$\displaystyle\sum$\cr\hidewidth$\displaystyle\int$\hidewidth\cr}}
  {\ooalign{\raisebox{.14\height}{\scalebox{.7}{$\textstyle\sum$}}\cr\hidewidth$\textstyle\int$\hidewidth\cr}}
  {\ooalign{\raisebox{.2\height}{\scalebox{.6}{$\scriptstyle\sum$}}\cr$\scriptstyle\int$\cr}}
  {\ooalign{\raisebox{.2\height}{\scalebox{.6}{$\scriptstyle\sum$}}\cr$\scriptstyle\int$\cr}}
}

\graphicspath{{./}{./figs/}}


%\begin{figure}
%  \includegraphics[width=0.8\columnwidth]{filename}
%   \caption{fig caption}
%   \label{labelname}
%\end{figure}
  
\begin{document}

\title{Renormalization of 3D Luttinger action}
\author{Santanu Dey}
\date{\today}
\begin{abstract}
Three-dimensional Luttinger fermions have quadratic band touching dispersion with
Lifshitz-like dynamical scaling $z = 2$. We extend the single-particle Luttinger
Hamiltonian to include the effects of minimally coupled gauge fluctuations and
illustrate the long-wavelength behavior of the combined interacting system. Within
the parameter space of our model, we discover an analog of the non-Fermi liquid phase
found in Luttinger semimetals with static Coulomb interaction~\cite{moon13}. We further
show that dynamical gauge fluctuations can also lead to non-relativistic Lifshitz
scaling for the emergent photons in our theory. With an RG analysis, we describe the
stability and scaling properties of its fixed points in the presence of additional
cubic-symmetry allowed perturbations.
\end{abstract}

\maketitle

\section{The action and its conserved currents}

\subsection{The fermion sector}

Let us consider an U(1) gauge theory coupled to fermions with 
symmetry protected quadratic band touching. The action without the gauge fields 
has been discussed before \cite{murakami04,moon13}. In the rotation and particle-hole
symmetric limit, using $N_\Gamma=(d-1)(d/2+1)$ Clifford matrices $\{\Gamma^a\}$ 
of dimensions $4\times 4$, the Euclidean Lagrangian is expressed as
\cite{janssen15,boettcher16}:
\beq
&\L_{\rm L} = \psi^\dagger\left[\partial_0
+\frac{d_a(-i\partial)}{2m}\Gamma^a\right]\psi,\nonumber\\
&d_a\left(-i\partial\right) = -\sqrt{\frac{d}{2(d-1)}}
\Lambda_{a}^{ij}\partial_i\partial_j, \ 
\{\Gamma^a,\Gamma^b\} =2\delta_{ab},
\label{matter}
\eeq
where $\{\Lambda_a\}$ are the $(d-1)(d+2)/2$ real, symmetric 
$d\times d$ Gell-Mann matrices \cite{boettcher16,janssen15}. In three 
dimensional momentum space the band dispersions are 
\cite{boettcher16},
\beq
&d_1(k)=\frac{3}{2}\left(k_x^2-k_y^2\right), \
d_2(k)=\sqrt{3}k_x k_y,\nonumber\\
&d_3(k)=\sqrt{3}k_xk_z, \
d_4(k)=\sqrt 3 k_y k_z, \nonumber\\
&d_5(k)=\frac{1}{2}\left(2k_z^2-k_x^2-k_y^2\right)
\eeq
In our notation the greek indices span the 
space-time and the roman indices only span the space.
It can be shown that $\Gamma^{ab}=\frac{1}{2}[\Gamma^a,\Gamma^b]$ are the 
$N_\Gamma(N_\Gamma-1)/2$ generators 
of $SO(N_\Gamma)$ rotations \cite{zinn21}. Certain useful identities of the $\Gamma$ 
matrices are as follows:
\beq
\Gamma^b\Gamma^a\Gamma^b&=(2-N_\Gamma)\Gamma^a\nonumber\\
\Gamma^c\Gamma^a\Gamma^b\Gamma^c&=(N_\Gamma-4)\Gamma^a\Gamma^b+4\delta_{ab}\nonumber\\
\Gamma^d\Gamma^a\Gamma^b\Gamma^c\Gamma^d&=-2\Gamma^c\Gamma^b\Gamma^a-(N_\Gamma-4)\Gamma^a\Gamma^b\Gamma^c
\eeq
The symmetric Gell-mann matrices obey the identities \cite{murakami04,boettcher16,janssen15},
\beq
&\Lambda^{ij}_a\Lambda^{kl}_a
=\left(\delta_{ik}\delta_{jl}+\delta_{il}\delta_{jk}\right)
-\frac{2}{d}\delta_{ij}\delta_{kl},\nonumber\\
&\tr\left[\Lambda_a\Lambda_b\right]=2\delta_{ab},
\eeq
Using these identities it can be shown in a simple way 
that \cite{boettcher16}
\beq
&d_a(p)d_a(q)=\frac{d}{d-1}\left(\left(\vec p_\perp\cdot\vec q_\perp\right)^2-
\frac{1}{d}
p_\perp^2q_\perp^2\right),\nonumber\\
\implies&
d_a(q)d_a(q)=q_\perp^4.
\eeq

The global U(1) symmetry of the action \eqref{matter} is manifest with the global transformations 
$\psi\mapsto e^{ie\alpha}\psi$. The corresponding global currents are given by,
\beq
j_0&=ie\psi^\dagger\psi,\nonumber\\
j_i&=ie\sqrt{\frac{d}{2(d-1)}}\frac{\Lambda_a^{ij}}{2m}
\left[
  \psi^\dagger\Gamma^a\partial_j\psi
  -
  \partial_j\psi^\dagger\Gamma^a\psi
\right],
\label{global_current}
\eeq
with the current conservation condition,
\beq
\partial_\mu j_\mu=0,
\eeq
dictated by the symmetry.
We seek to elevate the global U(1) symmetry of the model to a local U(1) gauge symmetry
by introducing vector gauge fields $A_\mu$ which are minimally coupled with the action.
The inclusion of the minimal coupling is achieved by replacing 
all the normal derivatives of the action with covariant derivatives
$D_\mu=\partial_\mu+ieA_\mu$. With the covariant derivatives, the action reads as:
\beq
\L'_{\rm L} &= \psi^\dagger\left[\partial_0+ieA_0
-\sqrt{\frac{d}{2(d-1)}}\Lambda_a^{ij}\Gamma^a\frac{D_i D_j}{2m}\right]\psi,\nonumber\\
\label{cov_matter}
\eeq
with the rule that under the local U(1) transformations both the fermionic fields and
the vector gauge fields transform simultaneously with the formula,
\beq
 {\rm U(1)}&:\psi\mapsto e^{ie\varphi(x)}\psi,\nonumber\\
{\rm U(1)}&:A_\mu(x)\mapsto A_\mu(x) -\partial_\mu\varphi(x),\\
\label{gauge}
\eeq
where $x=(\tau,\vec x)$ denotes the space-time coordinate. Similarly for the frequency-momentum
space we shall use the notation $k=(k_0,\vec k)$. 
It is straightforward to check that under the gauge transformations [Eq.~\eqref{gauge}]
the action of the covariant derivative changes as,
\beq
{\rm U(1)}&:D_\mu\psi\mapsto (\partial_\mu+i e A_\mu(x)-ie\partial_\mu\varphi(x))
e^{ie\varphi(x)}\psi
=e^{ie\varphi(x)}D_\mu\psi,
\eeq
which in turn ensures
that the minimally coupled action written down above \eqref{cov_matter} 
remains invariant under the gauge transformation.

\subsection{The photon sector}

With the fermionic sector well defined we need a gauge invariant description of the 
gauge field action. We are interested in an U(1) gauge theory in (3+1) dimension. The gauge 
invariant Euclidean Maxwell action is given by \cite{srednicki07,zinn21}
\beq
\L_{\rm Maxwell}=\frac{1}{4}F^2_{\mu\nu}+j^\mu A_\mu,
\eeq
with the matter current $j^\mu$ defined above \eqref{global_current}. 
However, we are in a situation where the Lorentz invariance is absent in the fermionic sector, 
and naturally, space and time has different engineering 
dimensions. This scenario has been explored in the context of the pure 
U(1) gauge theory of 
pyrochlore compounds \cite{moessner03,hermele04}. From the fermion sector it is imperative 
that before any loop correction the engineering dimension of space and time is 
respectively $[x^{-1}]=1$ and $[\tau^{-1}]=z_0=2$. The appropriate 
U(1) gauge theory in this situation 
goes by the name of $z_0=2$ Lifshitz electrodynamics \cite{das10,alexandre13} and its 
properties have been studied in details \cite{das10,alexandre13,anselmi10,farias12}. 
The Lagrangian is most compactly written in Euclidean space-time as,
\beq
\L_{\rm U(1)} &= 
\frac{1}{2}
F_{0i}F_{0i}+
\frac{1}{4}
F_{ij}
\left(u^2+v^2
\left[-
  \nabla^2\right]\right)
F_{ij}
+ j_\mu A_\mu, \text{ with }
F_{\mu\nu} = D_\mu A_\nu - D_\nu A_\mu.
\label{gauge_act}
\eeq
As a consequence of gauge invariance, the photon propagator obtained from this action is singular.
To determine the photon propagator we need
to work on a fixed gauge. In a Lorentz invariant theory a gauge fixing function 
$G(A)=\partial_\mu A_\mu$ can be set to zero because under a gauge transformation 
$A_\mu(x)\mapsto A^\alpha_\mu(x)=A_\mu(x)+\partial_\mu\alpha(x)$ the function transforms as
\beq
G(A^\alpha_\mu(x))=\partial_\mu A_\mu(x)+\partial^2\alpha(x),
\eeq
and we can always choose and $\alpha(x)$ such that the function vanishes. The gauge fixing
function can in principle be completely arbitrary. In the current
case a suitable 
gauge fixing function for the action is given by \cite{alexandre13} -
\beq
G(A_\mu,\omega)
=
\partial_0A_0(x)+L_g\partial_i A_i(x)-\omega(x)=0,
\label{gauge_orbit}
\eeq
where $\omega(x)$ is an arbitrary function and we have
defined a differential operator $L_g=u^2+v^2\left[-\nabla^2\right]$, 
In the path integral
we must only sum over the unique gauge configurations dictated by the 
gauge orbit $\omega(x)$ \eqref{gauge_orbit}. Various different gauge field
configurations which are related to each other through the 
$\alpha(x)$ reparameterization must be counted once. The original functional 
integral is given by
\beq
Z_{\rm U(1)}[j]=\int\mathcal D A e^{-\int_x \L_{\rm U(1)}(A_\mu,j_\mu)},
\eeq
where $\int_x=\int dx_0 d^dx_\perp$ is a short-hand notation for 
the $d+1$ dimensional space-time integral.
We wish to integrate over gauge field configurations which satisfy the 
gauge fixing condition $G(A_\mu,\omega)=0$. We shall therefore
sum over all gauge field reparameterization 
$A_\mu(x)\mapsto A^\alpha_\mu(x)$ that leaves the functional integral invariant
but also obey the gauge fixing condition. This can be carried out by the 
Fadeev-Popov prescription where one uses the identity \cite{peskin95}, 
\beq
1=\int\mathcal D\alpha\det\left(\frac{\delta G(A_\mu^\alpha,\omega)}{\delta\alpha}
\right)
\delta\left(G(A_\mu^\alpha,\omega)\right),
\eeq
to modify the partition function of the gauge fields with a uniform unity weight,
\beq
Z_{\rm U(1)}[j]=\int\mathcal D\alpha\mathcal D A 
e^{-\int_x \L_{\rm U(1)}(A_\mu,j_\mu)}
\det\left(\frac{\delta G(A_\mu^\alpha,\omega)}{\delta\alpha}
\right)
\delta\left(G(A_\mu^\alpha,\omega)\right).
\eeq
With $\frac{\delta G(A_\mu^\alpha,\omega)}{\delta\alpha}
=\partial_0^2+L_g\partial_i^2
$ the determinant in this expression factors out. With the 
intention to make the photon propagator more simple, a non-local gauge orbit can be 
chosen for our Lorentz-violating gauge field action
\cite{anselmi10,farias12,alexandre13} with a gauge fixing function,
\beq
G(A_\mu,\varphi)
=\partial_0A_0(x)+L_g\partial_i A_i(x)
-L_g^{1/2}\varphi(x),
\eeq
where we shall sum over all functions $\varphi(x)$ with a Gaussian weight centered around
$\varphi(x)=0$. Finally, we make a gauge transformation 
$A_\mu\mapsto A_\mu+\partial_\mu\alpha$ in the functional integral over the gauge fields 
to obtain,
\beq
Z_{\rm U(1)}[j]
=&\int\mathcal D\varphi e^{-
\int_x\frac{\varphi(x)^2}{2\xi}}
\int\mathcal D\alpha
\det\left(\partial_0^2+L_g\partial_i^2\right)
\nonumber\\
&\times
\int \mathcal D A^\alpha
e^{-\int_x\mathcal L_{\rm U(1)}(A_\mu^\alpha,j_\mu)}
\delta\left(\partial_0 A^\alpha_0(x)+
L_g\partial_iA^\alpha_i(x)
-L_g^{1/2}\varphi(x)
\right)\nonumber\\
=&
\det\left(\partial_0^2+L_g\partial_i^2\right)
\int\mathcal D\alpha
\int \mathcal D A^\alpha
e^{-\int_x\mathcal L_{\rm U(1)}(A^\alpha_\mu,j_\mu)}
e^{-
\int_x\frac{\left(
L_g^{-1/2}
\partial_0 A^\alpha_0(x)+
L_g^{1/2}\partial_iA^\alpha_i(x)
\right)^2}{2\xi}}.
\eeq
It is clear in the second line that $\alpha(x)$ has become a dummy parameter and naturally
the functional integral
over $\alpha(x)$ simply factors out. The procedure therefore achieves its
goal of summing over unique gauge field configurations.
The Feynman-'t Hooft gauge $\xi=1$ further simplifies the gauge fixed action,
\beq
\mathcal L_{\rm gf}
&=\frac{1}{2}F_{0i}F_{0i}
+\frac{1}{4}F_{ij}L_gF_{ij}
+\frac{1}{2}\left(
L_g^{-1/2}\partial_0A_0
+L_g^{1/2}\partial_i A_i
\right)^2,
\eeq
and yields a simple non-singular gauge field propagator 
to \cite{anselmi10},
\beq
D^{\rm gf}_{\mu\nu}(k)&=\langle A_\mu(-k)A_\nu(k)\rangle
=\left(
\delta_{\mu0}\delta_{\nu0}
\frac{L_g(k_\perp^2)}{k_0^2+k_\perp^2L_g(k_\perp^2)}
+
\delta_{\mu i}\delta_{\nu j}
\frac{\delta_{ij}}{k_0^2+k_\perp^2L_g(k_\perp^2)}
\right),
\eeq
where $k_\perp^2=k_i^2$ is the length of the momentum $\vec k_\perp$. 
Henceforth, we shall use the gauge-fixed propagator 
$D^{\rm gf}$ without referring to the gauge fixing directly.

\section{RG analysis}

\subsection{The action}

The renormalization group analysis of the fully interacting problem 
begins by separating the action in the free and the interacting 
vertices part. For this analysis we must consider a gauge-invariant
Lagrangian
with symmetry allowed terms that are closed under renormalization.
For this reason, in addition to the terms discussed so far we also
include a Pauli coupling term,
\beq
\frac{\kappa}{2}\psi^\dagger \Gamma^{ab}\psi
\eta^i_{ab}\epsilon_{ijk}
F_{jk},
\eeq
with $\Gamma^{ab}=\frac{1}{2i}\left[\Gamma^a,\Gamma^b\right]$ and
$\eta^i_{ab}=-\eta^i_{ba}$ are the symbols that relate Spin 3/2 SU(2) 
matrices with $\Gamma_{ab}$ \cite{murakami04},
\beq
S^i&=\eta^i_{ab}\Gamma^{ab}, \nonumber  
\left[\eta^i,\eta^j\right]=-\frac{1}{4}\epsilon^{ij}_{k}\eta^k.
\eeq
We shall need some more properties of $\eta^i$ matrices from their
defining relation. For $D=4$ dimensional $\Gamma^{ab}$ matrices
we have the formula $\tr\left[\Gamma^{ab}\Gamma^{cd}\right]
=D\left(\delta_{ac}\delta_{bd}-\delta_{ad}\delta_{bc}\right)$, which
we may use to find,
\beq
\tr\left[S^i S^j\right]
&=-2D\eta^i_{ab}\eta^j_{ba},\nonumber\\
\tr\left[\eta^i\eta^j\right]
&=-\frac{5}{2D}\delta^{ij}
=-\frac{5}{8}\delta^{ij}.
\eeq
It also follows that $\eta^i\eta^i=c I$, because this operator commutes with
all the generators $\eta^i$. Using the trace identity we have,
\beq
\eta^i\eta^i=-\frac{15}{8N_\Gamma}I
\eeq
To relate $\eta^i_{ab}$ symbols with the $\Lambda_a^{ij}$
matrices we need the formula \cite{boettcher17},
\beq
\Gamma^a=\frac{1}{\sqrt{2d(d-2)}}\{S^i,S^j\}\Lambda_a^{ij},
\label{spin_formula}
\eeq
and use it in conjunction with
the $\eta$ and $\Lambda$ matrices. We find,
\beq
\eta^i_{ab}&=
\frac{1}{8}\Tr\left[S^i\Gamma^{ab}\right]
,\nonumber\\
&=\frac{1}{32d(d-2)i}
\Lambda_a^{mn}\Lambda_b^{rs}
\Tr\left[S^i\left[\{S^m,S^n\},\{S^r,S^s\}\right]\right]
\eeq
The $S^i$ commutators in the last line can be extracted by inverting 
Eq.~\eqref{spin_formula} above.
The Pauli coupling $\kappa$ is relevant in the infrared at $d=3$. 
Together, we 
have the action,
\beq
\mathcal L =& \mathcal L_{\rm gf} +\mathcal L_{\rm L} 
+ \mathcal L_{\rm int},\nonumber\\
\mathcal L_{B,\rm g} =& 
\frac{1}{2}F_{B,0i}F_{B,0i}
+\frac{1}{4}F_{B,ij}\left(u^2_B+v^2_B\left[-\nabla^2\right]\right)F_{B,ij}
,\nonumber\\
\mathcal L_{B,\rm L} =& \psi_B^\dagger\left[\partial_0
+\frac{\slashed d\left(-i\partial\right)}{2m_B}\right]\psi_B,\nonumber\\
\mathcal L_{B,\rm int} =& 
ie_B \Big[\psi_B^\dagger
  A_{B,0}\psi_B
  -
  \sqrt{\frac{d}{2(d-1)}}
  \frac{1}{2m_B}
  \left(
  \psi_B^\dagger\slashed\Lambda^{ij}A_{B,i}\partial_j\psi_B
  -
  \left(\partial_j\psi_B^\dagger\right)
  \slashed\Lambda^{ij}A_{B,i}\psi_B
  \right)
  \Big]\nonumber\\
&+\frac{e_B^2}{2m_B}
\psi_B^\dagger
\slashed d\left(A\right)
\psi_B+
\frac{\kappa_B}{2}\psi_B^\dagger \Gamma^{ab}\psi_B
\eta^i_{ab}\epsilon_{ijk}
\left(\partial_j A_{B,k}
-\partial_k A_{B, j}
\right),
\eeq
where $u^2_B>0, v^2_B>0$ for a positive definite (bounded) bosonic action,
and we have introduced the Dirac slash notation, $P_a\Gamma^a=\slashed P$.
The zero temperature Euclidean theory is embedded in (d+1) dimensions. 
The tree level engineering dimension of space and time can be set at 
$[x]=-1$ and $[\tau]=-z_0$. The engineering dimensions of the
various fields and couplings of the theory are as follows,
\beq
&\left[\psi\right] = d/2, \ \left[m^{-1}\right] = z_0-2, 
\ \left[A_0\right] = (d+z_0-2)/2, 
\ \left[A_i\right] = (d-z_0)/2, \nonumber\\
&\left[e^2\right] = 2+z_0-d,
\ \left[u^2\right] = 2z_0-2,
\ \left[v^2\right] = 2z_0-4,\nonumber\\
&\left[\kappa\right] =  (3z_0-d-2)/2
\eeq
With $z_0=2$ all the couplings of the action are dimensionless 
except from $e^2$ and $u$. 
It is clear that the theory is 
strongly coupled with the relevant coupling $u$ in any spatial dimensions
with $z_0=2$. On the other hand the choice of $z_0 = 1$ turns the 
fermion bands flat. To investigate whether the theory has Wilsoninan 
fixed point within its parameter space we manually tune the relevant
coupling to a critical value $u_c=0$ and study the effective theory 
within a critical hypersurface parameterized by the remaining
marginal couplings. 
The critical gauge field action therefore becomes,
\beq
\mathcal L^c_{\rm g}
=\frac{1}{2}F_{B,0i}F_{B,0i}+\frac{v^2_B}{4}F_{B,ij}
\left[-\nabla^2\right]F_{B,ij}
\eeq
However, the gauge charge $e$ is also strongly relevant in
the physical spatial dimension of interest $d=3$. 
For perturbative control on the effective critical theory we follow Moon et al. 
\cite{moon13} to work in $d=4-\epsilon$ 
dimensions such that all the couplings of the theory are marginally relevant or 
irrelevant except from $u$ which we critically tune to zero. 
Along with 
the compact temporal integral (Euclidean time-cycle) we are therefore undertaking 
a $5-\epsilon$
expansion of the theory.
We start by renormalizing the bare fields and parameters as,
\beq
&\psi_B=\sqrt{Z_\psi}\psi, \ A_{B,\mu} = \sqrt{Z_A} A_\mu, \nonumber\\
&m^{-1}_B= Z_{m^{-1}} m^{-1}, \ e_B\mapsto Z_e e, \
v_B\mapsto Z_v v, \nonumber\\
&\kappa_B=Z_\kappa\kappa
\eeq
The renormalized Lagrangian is therefore given by,
\beq
\mathcal L_{\rm crit} 
=& \ \mathcal L^c_{\rm g} +\mathcal L_{\rm L} + \mathcal L_{\rm int},\nonumber\\
\mathcal L^c_{\rm g} =& \ 
\frac{Z_A}{2}F_{0i}F_{0i}+\frac{Z_A Z_v v^2}{4}F_{ij}
\left[-\nabla^2\right]F_{ij},
\nonumber\\
\mathcal L_{\rm L} = & \ 
Z_\psi\psi^\dagger\left[\partial_0
+Z_{m^{-1}}
\frac{\slashed d\left(-i\partial\right)}{2m}\right]\psi,\nonumber\\
\mathcal L_{\rm int} = & \ 
Z_e Z_\psi \sqrt{Z_A} i e \psi^\dagger\Big[
  A_0
  -Z_{m^{-1}}
  \sqrt{
    \frac{d}{2(d-1)}}
  \frac{\slashed\Lambda^{ij}
  \left(\partial_iA_j
  +A_j\partial_i
  +A_i\partial_j
  \right)}{2m}
  \Big]\psi\nonumber\\
&+Z_e^2Z_\psi Z_AZ_{m^{-1}} 
\frac{e^2}{2m}
\psi^\dagger
\slashed d\left(A\right)\psi
+Z_\kappa Z_\psi\sqrt Z_A
\frac{\kappa}{2}\psi^\dagger \Gamma^{ab}\psi
\eta^i_{ab}\epsilon_{ijk}
\left(\partial_j A_{k}
-\partial_k A_{j}
\right),
\eeq
Gauge invariance imposes further constraints on the renromalization parameters. 
We have,
\beq
&Z_\psi=Z_e Z_\psi\sqrt{Z_A} 
\text{ and } Z_\psi Z_{m^{-1}}=Z_e^2Z_\psi Z_AZ_{m^{-1}}
,\nonumber\\
\implies & Z_e\sqrt Z_A = 1.
\eeq
We redefine the remaining independent renormalization parameters as,
\beq
&Z_\psi=Z_1,\nonumber\\
&Z_\psi Z_{m^{-1}}=Z_2,\nonumber\\
&Z_A=Z_3,\nonumber\\
&Z_AZ_v=Z_4,\nonumber\\
&Z_\kappa Z_\psi\sqrt Z_A=Z_5.
\eeq
As a consequence it is possible to find the one-loop renormalization group coefficients
of the theory by studying its two-point functions alone.
Furthermore we shall expand around the tree level values of these renormalization 
factors, $Z_i=1+\delta Z_i=1+\delta_i$. This splits the Lagrangian 
between its tree level form and counter terms: 
\begin{align}
  \mathcal L =& \mathcal L^c_{\rm g} +\mathcal L_{\rm L} + \mathcal L_{\rm int}
  +\mathcal L_{\rm ct},\nonumber\\
  \mathcal L^c_{\rm g} =& 
  \frac{1}{2}F_{0i}F_{0i}+\frac{v^2}{4}F_{ij}
  \left[-\nabla^2\right]F_{ij},\nonumber\\
  \mathcal L_{\rm L} =& \psi^\dagger\left[\partial_0
  +\frac{\slashed d(-i\partial)}{2m}\right]\psi,\nonumber\\
  \mathcal L_{\rm int} =& 
  ie \Big[\psi^\dagger
  A_0\psi
  -
  \sqrt{\frac{d}{2(d-1)}}
  \frac{1}{2m}
  \left(
  \psi^\dagger\slashed\Lambda^{ij}A_i\partial_j\psi
  -
  \left(\partial_j\psi^\dagger\right)
  \slashed\Lambda^{ij}A_i\psi
  \right)
  \Big]+\frac{e^2}{2m}\psi^\dagger
  \slashed d(A)\psi\nonumber\\
  &
  +\kappa \psi^\dagger
  \sigma_{ij}
  \left(\partial_i A_j
  -\partial_j A_i
  \right)\psi
  ,\nonumber\\
  \mathcal L_{\rm ct}
  =&
  \delta_3\frac{1}{2}F_{0i}F_{0i}
  +\delta_4\frac{v^2}{2}F_{ij}F_{ij}
  \nonumber\\
  &
  +\delta_1\psi^\dagger\left[\partial_0+ieA_0\right]\psi
  +\delta_2\psi^\dagger\left[
    \frac{
      \slashed{d}(-i\partial+e A)
    }{2m}\right]\psi
    +
  \delta_5\frac{\kappa}{2}\psi^\dagger \Gamma^{ab}\psi
\eta^i_{ab}\epsilon_{ijk}
\left(\partial_j A_{k}
-\partial_k A_{j}
\right)
\end{align}
We have all the pieces in hand to evaluate the renormalized perturbation 
series diagrammatically in the momentum space. We adopt the following
convention for the fourier transformation of the fields,
\beq
A_\mu(x)&=\int\frac{d^{d+1}k}{(2\pi)^{d+1}}A_\mu(k)e^{ik.x},\\
\psi(x)&=\int\frac{d^{d+1}k}{(2\pi)^{d+1}}\psi(k)e^{ik.x}.\\
\eeq
As a consistent choice of notation, momentum entering a feynman diagram
vertex shall be considered positive, while the momentum leaving a 
vertex shall be considered negative. Defining a function,
$
\slashed d^i(k)=\sqrt{\frac{d}{2(d-1)}}\slashed\Lambda^{ij}k_j,
$
we can compactly write down the momentum space
ferynman rules for the action:
\begin{fmffile}{feynman-rules}
    \begin{align}
        \begin{gathered}
          \fmfframe(4,4)(4,0){
            \begin{fmfgraph*}(20,10)
                \fmfleft{v1}
                \fmfright{v2}
                \fmflabel{$\alpha$}{v1}
                \fmflabel{$\beta$}{v2}
                \fmf{fermion,label=$k$}{v1,v2}  
            \end{fmfgraph*}
          }
        \end{gathered}=& \
        \langle\psi_\alpha(k)\psi_\beta^\dagger(k)
        \rangle = G_{\alpha\beta}(k)
        =\left[\frac{-ik_0+\slashed{d}(k)/(2m)} 
        {k_0^2+
        k^4_\perp/(4m^2)}\right]_{\alpha\beta} 
        \nonumber\\
        \begin{gathered}
          \fmfframe(4,4)(4,0){
          \begin{fmfgraph*}(20,10)
                \fmfleft{v1}
                \fmfright{v2}
                \fmflabel{$\mu$}{v1}
                \fmflabel{$\nu$}{v2}
                \fmf{photon,label=$k\rightarrow$}{v1,v2}  
            \end{fmfgraph*}
          }
        \end{gathered}
        =& \ 
        \langle A_\mu(-k)A_\nu(k)\rangle=
        D^{\rm gf}_{\mu\nu}(k)
        =\left(
        \delta_{\mu0}\delta_{\nu0}
        \frac{v^2 k_\perp^2}{k_0^2+v^2 k_\perp^4}
        +
        \delta_{\mu i}\delta_{\nu j}
        \frac{\delta_{ij}}{k_0^2+v^2 k_\perp^4}
        \right),
        \nonumber\\
        \begin{gathered}
          \fmfframe(8,8)(4,0){
            \begin{fmfgraph*}(20,20)
                \fmfleft{i1,i2}
                \fmfright{o1}
                \fmf{fermion,label=$k+q$}{i1,v1}  
                \fmf{fermion,label=$k$}{v1,i2}
                \fmf{photon,label=$q \rightarrow $}{o1,v1}  
                \fmflabel{$\alpha$}{i1}
                \fmflabel{$\beta$}{i2}
                \fmflabel{$\mu$}{o1}
            \end{fmfgraph*}
          }
        \end{gathered}
        =& \ C^{\mu}_{\alpha\beta}
        (k+q,-q,-k)
        = -e\Big[i\delta_{\mu 0}+
        \frac{\left(\slashed d^i(k+q)-\slashed d^i(-k)\right)}{2m}
        \delta_{\mu i}
        \Big]_{\alpha\beta}\nonumber\\
        &
        -\kappa
        \left[\Gamma^{ab}\right]_{\alpha\beta}
        \eta^i_{ab}\epsilon_{ijk}
        \left(-i q_j\right)\delta_{\mu k},\nonumber\\
        \begin{gathered}
          \fmfframe(8,8)(4,0){
            \begin{fmfgraph*}(20,20)
                \fmfleft{i1,i2}
                \fmfright{o1,o2}
                \fmf{fermion,label=$k_1$}{i1,v1}  
                \fmf{fermion,label=$k_1+k_2-k_3$}{v1,i2}  
                \fmf{photon,label=$\nwarrow k_2$,label.side=right,label.dist=-2}{o1,v1}  
                \fmf{photon,label=$k_3\nearrow$,label.side=left,label.dist=-2}{v1,o2}  
                \fmflabel{$\alpha$}{i1}
                \fmflabel{$\beta$}{i2}
                \fmflabel{$\mu$}{o1}
                \fmflabel{$\nu$}{o2}
            \end{fmfgraph*}
          }
        \end{gathered}
        =& \ V^{\mu\nu}_{\alpha\beta}(k_1,k_2,-k_3,-k_1-k_2+k_3) 
        =-\frac{e^2}{m}\sqrt{\frac{d}{2(d-1)}} 
        \delta_{\mu i}\delta_{\nu j}
        \left[\slashed\Lambda^{ij}\right]_{\alpha\beta},
    \end{align}
\end{fmffile}

\noindent
where $\{u^\mu\}$ are the Euclidean unit vectors and the factor of 2 in the quartic vertex 
comes from the symmetry permutation of the 
photon fields. The Feynman diagram of the counter term vertices are as follows:
\begin{fmffile}{counter_term}
    \begin{align}
        \begin{gathered}
          \fmfframe(4,4)(4,0){
            \begin{fmfgraph*}(20,10)
                \fmfleft{i}
                \fmfright{o}
                \fmflabel{$\alpha$}{i}
                \fmflabel{$\beta$}{o}
                \fmf{fermion,label=$k$}{i,v}  
                \fmf{fermion,label=$k$}{v,o} 
                \fmfv{decor.shape=circle,decor.filled=empty}{v}
            \end{fmfgraph*}
          }
        \end{gathered}=& \
        \delta_{\Sigma}=
        -\left[
          \delta_1\left(i k_0\right)
        +
        \delta_2\left(
        \frac{\slashed d(k)}{2m}
        \right)
        \right]_{\alpha\beta},
        \nonumber\\
        \begin{gathered}
          \fmfframe(4,4)(4,0){
            \begin{fmfgraph*}(20,10)
                \fmfleft{i}
                \fmfright{o}
                \fmflabel{$\mu$}{i}
                \fmflabel{$\nu$}{o}
                \fmf{photon,label=$k\rightarrow$}{i,v}  
                \fmf{photon,label=$k\rightarrow$}{v,o} 
                \fmfv{decor.shape=square,decor.filled=empty}{v}
            \end{fmfgraph*}
          }
        \end{gathered}=& \
        \delta_{\Pi} =
        -\frac{\delta_3}{2}
        \left\{k^2\delta_{\mu\nu} -k_\mu k_\nu-
        \delta_{\mu i}\delta_{\nu j}
        \left(
        k_\perp^2\delta_{ij} -k_i k_j
        \right)
        \right\}
        \nonumber\\
        & \ -
        \delta_4  
        \delta_{\mu i}\delta_{\nu j}
        \frac{v^2
        k_\perp^2}{2}\left(\delta_{ij}k_\perp^2-k_ik_j\right)
        ,\nonumber\\
        \begin{gathered}
          \fmfframe(8,8)(4,0){
            \begin{fmfgraph*}(20,20)
                \fmfleft{i1,i2}
                \fmfright{o1}
                \fmf{fermion,label=$k+q$}{i1,v1}  
                \fmf{fermion,label=$k$}{v1,i2}
                \fmf{photon,label=$q \rightarrow$}{o1,v1}  
                \fmfv{decor.shape=triangle,decor.filled=empty}{v1}
                \fmflabel{$\alpha$}{i1}
                \fmflabel{$\beta$}{i2}
                \fmflabel{$\mu$}{o1}
            \end{fmfgraph*}
          }
        \end{gathered}
        =& \ \delta_C = -e \Big[
          \delta_1\left(
          i\delta_{\mu 0}\right)+
        +
        \delta_2\left(
        \frac{\left(\slashed d^i(k+q)-\slashed d^i(-k)\right)}{2m}
        \right)\delta_{\mu i}
        \Big]_{\alpha\beta}\nonumber\\
      &        -\delta_5\kappa
        \left[\Gamma^{ab}\right]_{\alpha\beta}
        \eta^i_{ab}\epsilon_{ijk}
        \left(-i q_j\right)\delta_{\mu k}      ,\nonumber\\
        \begin{gathered}
          \fmfframe(8,8)(4,0){
            \begin{fmfgraph*}(20,20)
                \fmfleft{i1,i2}
                \fmfright{o1,o2}
                \fmf{fermion,label=$k_1$}{i1,v1}  
                \fmf{fermion,label=$k_1+k_2-k_3$,label.side=left,label.dist=-2}{v1,i2}  
                \fmf{photon,label=$\nwarrow k_2$,label.side=right,label.dist=-2}{o1,v1}  
                \fmf{photon,label=$k_3\nearrow$,label.side=left,label.dist=-2}{v1,o2}  
                \fmfv{decor.shape=diamond,decor.filled=empty}{v1}
                \fmflabel{$\alpha$}{i1}
                \fmflabel{$\beta$}{i2}
                \fmflabel{$\mu$}{o1}
                \fmflabel{$\nu$}{o2}
            \end{fmfgraph*}
          }
        \end{gathered}
        =& \ 
        \delta_V=
        -\delta_2\frac{e^2}{m}
        \sqrt{\frac{d}{2(d-1)}}
        \delta_{\mu i}\delta_{\nu j}
        \left[
          \slashed\Lambda^{ij}
        \right]_{\alpha\beta}
    \end{align}
\end{fmffile}

\noindent
It is clear that the counter terms are determined from the self-energy correction
to the inverse propagator alone. The full 
renormalized perturbation series can now be setup in terms of these diagrams. 
Following standard technique we sum up all one-particle irreducible 
(1PI) diagrams \cite{peskin95} and find,
\begin{fmffile}{pert_theory}
    \begin{align}
        \begin{gathered}
          \fmfframe(4,4)(4,0){
            \begin{fmfgraph*}(20,10)
                \fmfleft{v1}
                \fmfright{v2}
                \fmf{dbl_plain_arrow}{v1,v2}  
            \end{fmfgraph*}
          }
        \end{gathered}=& \
        \begin{gathered}
          \fmfframe(4,4)(4,0){
            \begin{fmfgraph*}(20,10)
                \fmfleft{v1}
                \fmfright{v2}
                \fmf{fermion}{v1,v2}  
            \end{fmfgraph*}
          }
        \end{gathered}
        +
        \begin{gathered}
          \fmfframe(4,4)(4,0){
            \begin{fmfgraph*}(40,10)
                \fmfleft{i}
                \fmfright{o}
                \fmf{fermion}{i,v1}  
                \fmf{dbl_plain_arrow}{v2,o}
                \fmf{plain,left,tension=0.2}{v1,v2,v1}
                \fmf{phantom,label.dist=0,label=$\Sigma^{\rm 1PI}_\psi$}{v1,v2}
            \end{fmfgraph*}
          }
        \end{gathered}
        \nonumber\\
        \nonumber\\
        \nonumber\\
        \Sigma_\psi^{1PI}
        =& \
        \begin{gathered}
          \fmfframe(4,4)(4,0){
            \begin{fmfgraph*}(20,10)
                \fmfleft{v1}
                \fmfright{v2}
                \fmf{fermion}{v1,v2}  
                \fmf{photon,left,tension=0.2}{v1,v2}  
            \end{fmfgraph*}
          }
        \end{gathered}
        +
        \begin{gathered}
          \fmfframe(4,4)(4,0){
            \begin{fmfgraph*}(20,10)
                \fmfleft{i}
                \fmfright{o}
                \fmf{phantom}{i,v,o}  
                \fmf{photon,right,tension=0.5}{v,v}  
            \end{fmfgraph*}
          }
        \end{gathered}+
        \begin{gathered}
          \fmfframe(4,4)(4,0){
            \begin{fmfgraph*}(20,10)
                \fmfleft{i}
                \fmfright{o}
                \fmf{phantom}{i,v}  
                \fmf{phantom}{v,o} 
                \fmfv{decor.shape=circle,decor.filled=empty}{v}
            \end{fmfgraph*}
          }
        \end{gathered}+\dots
        \nonumber\\
        \begin{gathered}
          \fmfframe(4,4)(4,0){
            \begin{fmfgraph*}(20,10)
                \fmfleft{v1}
                \fmfright{v2}
                \fmf{dbl_wiggly}{v1,v2}  
            \end{fmfgraph*}
          }
        \end{gathered}=& \
        \begin{gathered}
          \fmfframe(4,4)(4,0){
            \begin{fmfgraph*}(20,10)
                \fmfleft{v1}
                \fmfright{v2}
                \fmf{photon}{v1,v2}  
            \end{fmfgraph*}
          }
        \end{gathered}
        +
        \begin{gathered}
          \fmfframe(4,4)(4,0){
            \begin{fmfgraph*}(40,10)
                \fmfleft{i}
                \fmfright{o}
                \fmf{photon}{i,v1}  
                \fmf{dbl_wiggly}{v2,o}
                \fmf{plain,left,tension=0.2}{v1,v2,v1}
                \fmf{phantom,label.dist=0,label=$\Sigma^{\rm 1PI}_A$}{v1,v2}
            \end{fmfgraph*}
          }
        \end{gathered}
        \nonumber\\
        \nonumber\\
        \nonumber\\
        \Sigma_A^{1PI}
        =& \
        \begin{gathered}
          \fmfframe(4,4)(4,0){
            \begin{fmfgraph*}(20,10)
                \fmfleft{v1}
                \fmfright{v2}
                \fmf{fermion,left,tension=0.2}{v1,v2,v1}  
            \end{fmfgraph*}
          }
        \end{gathered}
        +
        \begin{gathered}
          \fmfframe(4,4)(4,0){
            \begin{fmfgraph*}(20,10)
                \fmfleft{i}
                \fmfright{o}
                \fmf{phantom}{i,v,o}  
                \fmf{fermion,right,tension=0.5}{v,v}  
            \end{fmfgraph*}
          }
        \end{gathered}
        +\begin{gathered}
          \fmfframe(4,4)(4,0){
            \begin{fmfgraph*}(20,10)
                \fmfleft{i}
                \fmfright{o}
                \fmf{phantom}{i,v}  
                \fmf{phantom}{v,o} 
                \fmfv{decor.shape=square,decor.filled=empty}{v}
            \end{fmfgraph*}
          }
        \end{gathered}+\dots
        \nonumber\\
    \end{align}
\end{fmffile}

\noindent
In terms of the vertex functions and the propagators the 
diagrams are evaluated to
\beq
\Sigma^{\rm 1L}_\psi(k)
=& 
\int_q C^{\mu}(k+q,-q,-k)G(k+q)
C^{\nu}(k,q,-k-q)D^{\rm gf}_{\mu\nu}(q)\nonumber\\
&+
\int_q V^{\mu\nu}(k,q,-q,-k)D^{\rm gf}_{\mu\nu}(q)
+\delta_{\Sigma},\nonumber\\
\Sigma_A^{\mu\nu,\rm 1L}(k) 
=&-
\int_q \tr\left[
  C^\mu(q+k,-k,q)G(q+k)
  C^\nu(q,k,-q-k)G(q)
  \right]\nonumber\\
&-
\int_q \tr\left[V^{\mu\nu}(q,k,-k,-q)G(q)\right]
+\delta_\Pi,
\eeq
where $\int_q=\int\frac{dq_0}{2\pi}\frac{d^dq_\perp}{(2\pi)^d}$ is the shorthand 
for the $d+1$ dimensional integral over internal frequency-momenta $q$.
To determine the renormalization parameters we shall demand the finiteness of the 
tree level action with respect to an artificial renormalization scale 
$\mu$. Which implies that the 
divergent pieces in the self energy expressions should
identically cancel. This criteria shall determine the divergent counter terms 
of the action.

It can be shown that the only external momentum dependence of 
self-energies are through the set of functions $(k_0,d_a(k),d_a^i(k))$. Therefore, we
may perform a Taylor expansion to obtain the leading order external momentum
dependence of the self-energies in terms of this function. We can introduce some
additional parameters $t,u,v$ and obtain,
\beq
\Sigma(tk_0,u d_a(k))
\eeq

\subsection{Useful formulae for the one-loop integrals}

We shall compute the frequency part of the loop integrals exactly as they involve no
divergences. For the remaining $d$ dimensional momentum integral we shall first
resolve the complicated denominators of the frequency-integrated 
loop expressions using 
the Feynman's parametric integral formulae\cite{kleinert01}:
\beq
\frac{1}{A^aB^b}&=\frac{\Gamma(a+b)}{\Gamma(a)\Gamma(b)}
\int^1_0dx\frac{x^{a-1}(1-x)^{b-1}}{[Ax+B(1-x)]^{a+b}},\nonumber\\
\frac{1}{A_1^{a_1}\dots A_n^{a_n}}
&=\frac{\Gamma(a_1+\dots+a_n)}{\Gamma(a_1)\dots\Gamma(a_n)}
\int^1_0 dx_1\dots\int^1_0dx_n\frac{\delta(1-(x_1+\dots+x_n))x_1^{a_1-1}\dots x_n^{a_n-1}}
{(x_1A_1+\dots+x_nA_n)^{a_1+\dots+a_n}}.
\eeq
Following the Feynman parameterization, a momentum shift puts the generic loop integral
in the form,
$
\int_{q_\perp} f_1(q_\perp^2)f_2(q_i,k_\mu)
$,
which can be further reduced by considering the integral identity,
\beq
\int\frac{d^dq}{(2\pi)^d}f(q^2)q_iq_j
=\frac{\delta_{ij}}{d}
\int\frac{d^dq}{(2\pi)^d}f(q^2)q^2.
\eeq
In fact we shall need similar identities involving multiple products of internal
momentum with open indices. In such types of integral 
the only non-zero contribution comes from cases where all the internal momentum
indices are pairwise contracted. It is straightforward 
to show this with the help of the Laplace transform of the (well-behaved) 
function $f(q^2)$. With the transform the integrals of the type,
\beq
K_{ijkl}&=\int\frac{d^dq}{(2\pi)^d}f(q^2)q_i q_j q_k q_l,\\
&=\int^\infty_0dt\int\frac{d^dq}{(2\pi)^d}
\tilde f(t)e^{-q^2t}q_iq_jq_kq_l,
\eeq
can be recast in terms of a generating variable $Y_i$ such that,
\beq
K_{ijkl}
&=
\left.
\frac{\partial^4}{\partial Y_i
\partial Y_j
\partial Y_k
\partial Y_l
}
\left[
\int^\infty_0dt\int\frac{d^dq}{(2\pi)^d}
\tilde f(t)e^{-q^2t+Y_iq_i}
\right]
\right|_{Y_i\rightarrow 0}
,\nonumber\\
&=
\left.
\frac{\partial^4}{
  \partial Y_i
  \partial Y_j
  \partial Y_k
  \partial Y_l
}
\left[
\int^\infty_0dt\tilde f(t)
\frac{
  e^{Y_i^2/(4t)}
}{(2\pi)^{d/2}(2t)^{d/2}}
\right]
\right|_{Y_i\rightarrow 0}
,\nonumber\\
&=
\int^\infty_0dt
\frac{
\tilde f(t)
}{(2\pi)^{d/2}(2t)^{d/2}}
\left(
\frac{\delta_{ij}}{2t}
\frac{\delta_{kl}}{2t}
+
\frac{\delta_{ik}}{2t}
\frac{\delta_{jl}}{2t}
+
\frac{\delta_{il}}{2t}
\frac{\delta_{jk}}{2t}
\right).
\eeq
On the other hand we find that,
\beq
\int\frac{d^dq}{(2\pi)^d}f(q^2)q^4
&=\int^\infty_0 dt\int\frac{d^dq}{(2\pi)^d}
\tilde f(t)e^{-q^2t}q^4,\nonumber\\
&=\int^\infty_0 dt
\tilde f(t)
\frac{\partial^2}{\partial^2 t}
\int\frac{d^dq}{(2\pi)^d}e^{-q^2t},\nonumber\\
&=\int^\infty_0 dt
\tilde f(t)
\frac{\partial^2}{\partial^2 t}
\left(
\frac{1}{(2\pi)^{d/2}(2t)^{d/2}}\right),\nonumber\\
&=\int^\infty_0 dt
\frac{\tilde f(t)}{(2\pi)^{d/2}(2t)^{d/2}}
\left(\frac{d(d+2)}{4t^2}\right).\nonumber\\
\eeq
Combining these two results we find that,
\beq
K_{ijkl}&=\int\frac{d^dq}{(2\pi)^d}f(q^2)q_i q_j q_k q_l
=\frac{1}{d(d+2)}
\int\frac{d^dq}{(2\pi)^d}f(q^2)q^4
\left(
\delta_{ij}
\delta_{kl}
+
\delta_{ik}
\delta_{jl}
+
\delta_{il}
\delta_{jk}
\right),
\eeq
and similarly,
\beq
K_{ijklmn}&=\int\frac{d^dq}{(2\pi)^d}f(q^2)q_i q_j q_k q_l q_m q_n
=\frac{1}{d(d+2)(d+4)}
\int\frac{d^dq}{(2\pi)^d}f(q^2)q^4
\left(
\delta_{ij}
\delta_{kl}\delta_{mn}
+
\dots
\right).
\eeq

\subsection{The fermion self-energy}

We start with the one-loop renormalized fermion self-energy 
$\Sigma^{\rm 1L}_\psi(k)$. The full expression is,
\beq
\Sigma^{\rm 1L}_\psi(k)
=&\int_q
\Bigg[
e\left[i\delta_{\mu 0}+
        \frac{
          \slashed d^i(q+2k)
        }{2m}
        \delta_{\mu i}
        \right]\nonumber\\
& \ \times
        \left[\frac{-i(k_0+q_0)
        +\slashed d(k+q)/(2m)
        } 
        {(k_0+q_0)^2+
        d(k+q)^2/(4m^2)}\right]\nonumber\\
& \ \times
e\left[i\delta_{\nu 0}+
        \frac{
          \slashed d^i(q+2k)
        }{2m}
        \delta_{\nu i}
        \right]\nonumber\\
& \ \times
    \left(
        \delta_{\mu0}\delta_{\nu0}
        \frac{v^2 q_\perp^2}{q_0^2+v^2 q_\perp^4}
        +
        \delta_{\mu i}\delta_{\nu j}
        \frac{\delta_{ij}}{q_0^2+v^2 q_\perp^4}
        \right)
\Bigg]
  -\left[
    \delta_1\left(ik_0\right)
    +\delta_2\left(
    \frac{\slashed d(k)}{2m}\right)
    \right].
\eeq
The contribution from the quartic vertex identically vanishes in this case. 
We shall consider the case for non-zero external frequency and momenta separately.
For the $k=(k_0,0)$ we obtain,
\beq
\tr\left[\Sigma^{\rm 1L}_\psi(k_0,0)\right]
=&\int_q
\tr\Bigg[
e\left[i\delta_{\mu 0}+
        \frac{
          \slashed d^i(q)
        }{2m}
        \delta_{\mu i}
        \right]\left[\frac{-i (q_0+k_0)
        +\slashed d(q)/(2m)
        } 
        {(q_0+k_0)^2+
        d(q)^2/(4m^2)}\right]\nonumber\\
& \ \times
e\left[i\delta_{\nu 0}+
        \frac{
          \slashed d^i(q)
        }{2m}
        \delta_{\nu i}
        \right]\left(
        \delta_{\mu0}\delta_{\nu0}
        \frac{v^2 q_\perp^2}{q_0^2+v^2 q_\perp^4}
        +
        \delta_{\mu i}\delta_{\nu j}
        \frac{\delta_{ij}}{q_0^2+v^2 q_\perp^4}
        \right)
\Bigg]
  -4\delta_1\left(ik_0\right),\nonumber\\
  =&-\frac{2i e^2(d-8m^2v^2+2)}{4m^2}
  \int_q\frac{(q_0+k_0)q^2}{(q_0^2+v^2q^4)((q_0+k_0)^2+\frac{q^4}{4m^2})}
  -4\delta_1(ik_0)
\eeq

\subsection{The photon self-energy}

Now we turn to the counter terms in the photon sector. It shall be convenient to introduce 
the transverse and longitudinal projectors associated with the photon propagator, defined
respectively as,
\beq
P^T_{\mu\nu}(k)=\delta_{\mu\nu}-\frac{k_\mu k_\nu}{k^2},
\ 
P^L_{\mu\nu}(k)=\frac{k_\mu k_\nu}{k^2}.
\eeq
We shall also define similar but separate 
projectors that act on the spatial components alone,
\beq
p^T_{ij}(k)=\delta_{ij}-\frac{k_i k_j}{k_\perp^2},
\ 
p^L_{ij}(k)=\frac{k_i k_j}{k_\perp^2}.
\eeq
With the aide of these projectors the
one-loop renormalized photon self-energy can be written down as,
\begin{align}
\Sigma_A^{\mu\nu,\rm 1L}(k_0,\vec k_\perp) 
=&-
\int_q
\tr\Bigg[
  e\left[
    i\delta_{\mu0}
    +\frac{
      \delta_{\mu i}\slashed d^i(k+2q)
    }{2m}
    \right] \nonumber\\
  & \times
  \left[
    \frac{-i(q_0+k_0)+\frac{\slashed d(q+k)}{2m}}
    {\left(q_0+k_0\right)^2+
    \frac{|\vec q_\perp+\vec k_\perp|^4}{4m^2}
    }
    \right] \nonumber\\
  & \times
  e\left[
    i\delta_{\nu0}
    +\frac{
      \delta_{\nu i}\slashed d^i(k+2q)
    }{2m}
    \right]\left[
    \frac{-i q_0+\frac{\slashed d(q)}{2m}}
    {q_0^2+
    \frac{q_\perp^4}{4m^2}
    }
    \right] \Bigg]\nonumber\\
&-\delta_3
  \left(k^2P^T_{\mu\nu}(k)-k_\perp^2\delta_{\mu i}\delta_{\nu j}p^T_{ij}(k)\right)
  \nonumber\\
& \ -
\delta_4  
\delta_{\mu i}\delta_{\nu j}
  v^2
  k_\perp^4p^T_{ij}(k).
\end{align}
The contribution from the quartic vertex once again drops out identically due to the properties
of the $\Lambda_a$ matrices. The computation also greatly simplifies if we evaluate the 
self-energy at zero external frequency, $k_0 = 0$.

To extract the counter term $\delta_4$ we separate the spatial components of the photon self-energy
and evaluate it at the fixed external frequency $k_0 = 0$. We obtain,
\begin{align}
  &\Sigma_A^{ij,\rm 1L}(0,\vec k)\nonumber\\ 
  =&-
\int_q
\tr\Bigg[
  e\left[
    \frac{
      \slashed d^i(k+2q)
    }{2m}
    \right]\left[
    \frac{-iq_0+\frac{\slashed d(q+k)}{2m}}
    {q_0^2+
    \frac{|\vec q+\vec k|^4}{4m^2}
    }
    \right] \nonumber\\
  & \times
  e\left[
    \frac{
      \slashed d^j(k+2q)
    }{2m}
    \right]\left[
    \frac{-i q_0+\frac{\slashed d(q)}{2m}}
    {q_0^2+
    \frac{q_\perp^4}{4m^2}
    }
    \right] \Bigg]-\delta_4 v^2 k_\perp^4p^T_{ij}(k),\\
=&-\frac{e^2}{4m^2}
\int_q
  \frac{M_1^{ij}(\vec k,q_0,\vec q)}{
    \left(q_0^2+\frac{|\vec q+\vec k|^4}{4m^2}\right)
    \left(q_0^2+\frac{q_\perp^4}{4m^2}\right)
  }
  -\delta_4 v^2 k_\perp^4p^T_{ij}(k),\\
\end{align}
where,
\beq
M_1^{ij}(\vec k,q_0,\vec q)=\tr\left[
  \slashed d^i(k+2q)
  \left(-iq_0+\frac{\slashed d(q+k)}{2m}\right)
  \slashed d^j(k+2q)
  \left(-iq_0+\frac{\slashed d(q)}{2m}\right)
  \right],
\eeq
is the trace product in the numerator. We compute the frequency integral exactly using
residue theorem and obtain,
\beq
&\Sigma_A^{ij,\rm 1L}(0,\vec k)\nonumber\\ 
=&-\frac{e^2}{4m^2}
\int_{\vec q}
  \left(
  -\frac{1}{2}\frac{M_1^{ij}(\vec k,i |\vec k+\vec q|^2/(2m),\vec q)}
  {\frac{|\vec k+\vec q|^2}{2m}
  \left(
  \frac{|\vec k+\vec q|^2}{2m}
  -
  \frac{q_\perp^2}{2m}
  \right)
    \left(
  \frac{|\vec k+\vec q|^2}{2m}
  +
  \frac{q_\perp^2}{2m}
  \right)
  }
  +\frac{1}{2}\frac{M_1^{ij}(\vec k,i q_\perp^2/(2m),\vec q)}
  {\frac{q_\perp^2}{2m}
  \left(
  \frac{|\vec k+\vec q|^2}{2m}
  -
  \frac{q_\perp^2}{2m}
  \right)
    \left(
  \frac{|\vec k+\vec q|^2}{2m}
  +
  \frac{q_\perp^2}{2m}
  \right)
  }
  \right)\nonumber\\
  &-\delta_4 v^2 k_\perp^4p^T_{ij}(k),\nonumber\\
=&\frac{d e^2}{4(d-1)^2m}
\int_{\vec q}
\frac{M_2^{ij}(\vec k,\vec q)}{q_\perp^2
\left(q_\perp^2+2\vec q\cdot\vec k +k_\perp^2\right)
\left(q_\perp^2+\frac{2\vec q\cdot\vec k +k_\perp^2}{2}\right)
}\nonumber\\
=&\frac{2d e^2}{4(d-1)^2m}
\int_{\vec q}
\int^1_0 dx_1
\int^1_0 dx_2
\int^1_0 dx_3
\frac{\delta\left(1-\left(x_1+x_2+x_3\right)\right)
M_2^{ij}(\vec k,\vec q)}{\left(x_1 q_\perp^2+
x_2\left(q_\perp^2+2\vec q\cdot\vec k +k_\perp^2\right)
+x_3\left(q_\perp^2+\frac{2\vec q\cdot\vec k +k_\perp^2}{2}\right)\right)^3
},\nonumber\\
=&\frac{2d e^2}{4(d-1)^2m}
\int_{\vec q}
\int^1_0 dx_1
\int^1_0 dx_2
\int^1_0 dx_3 \nonumber\\
&\times \frac{\delta\left(1-\left(x_1+x_2+x_3\right)\right)
M_2^{ij}(\vec k,\vec q)}
{\left(q_\perp^2
+2\vec q\cdot\left(x_2\vec k+(x_3/2)\vec k\right)
+\left|x_2\vec k+(x_3/2)\vec k\right|^2
+\left(x_2+x_3/2\right)\left(1-
\left(x_2+x_3/2\right)
\right)k_\perp^2
\right)^3
},\nonumber\\
=&\frac{2d e^2}{4(d-1)^2m}
\int_{\vec q}
\int^1_0 dx_1
\int^1_0 dx_2
\int^1_0 dx_3 \frac{\delta\left(1-\left(x_1+x_2+x_3\right)\right)
M_2^{ij}(\vec k,\vec q-(x_2+x_3/2)\vec k)}
{\left(q_\perp^2+\left(x_2+x_3/2\right)\left(1-
\left(x_2+x_3/2\right)
\right)k_\perp^2
\right)^3
},\nonumber\\
\eeq
where in the last two lines we have used the delta function condition $x_1+x_2+x_3=1$ 
and shifted the internal momenta to simplify the integral.


\bibliography{luttinger_act}
\end{document}
